%% ============================================================
%% CHAPTER 6 — RQ3: TECHNO-ECONOMICS
%% ============================================================
\chapter{Techno-Economic Analysis (RQ3)}
\label{ch:results_rq3}

This chapter answers RQ3 on the fixed-source basis of
Section~\ref{sec:method_rq3}: each cycle is sized to the output it can
recover from the site stream (Table~\ref{tab:sitedesign}), costed at that
size, and evaluated at the site electricity price.
Section~\ref{sec:rq3_deterministic} gives the deterministic results,
Section~\ref{sec:rq3_recovery} examines why the ranking differs from the
efficiency ranking, Section~\ref{sec:rq3_montecarlo} the Monte Carlo
robustness and breakeven analysis.

\section{Deterministic LCOE, NPV, Payback and IRR}
\label{sec:rq3_deterministic}

\begin{figure}[H]
\centering
\thesisfig{RQ3_LCOE_NPV_MONTE_CARLO.png}{1.0}{10cm}
\caption{Techno-economic results on a fixed-source basis (2026 AUD): net
power recovered from each stream, LCOE against the site electricity price,
project NPV, and Monte Carlo probability of positive NPV.}
\label{fig:rq3_results}
\end{figure}

\begin{table}[H]
\centering
\caption{Deterministic techno-economic results on a fixed-source basis
(7\,\% real discount rate, 20-year life, 85\,\% capacity factor, O\&M 2.5\,\%
of capital). Payback is discounted.}
\label{tab:rq3_results}
\renewcommand{\arraystretch}{1.25}
\footnotesize
\begin{tabularx}{\textwidth}{>{\raggedright\arraybackslash}p{0.9cm}>{\raggedright\arraybackslash}p{1.0cm}
  >{\centering\arraybackslash}X >{\centering\arraybackslash}X >{\centering\arraybackslash}X
  >{\centering\arraybackslash}X >{\centering\arraybackslash}X >{\centering\arraybackslash}X >{\centering\arraybackslash}X}
\toprule
\rowcolor{tableheader}
\textbf{Site} & \textbf{Cycle} & Output (kW) & TEC (kAUD) & TEC (AUD/kW) & LCOE (AUD/MWh) & NPV (kAUD) & Payback (yr) & IRR (\%) \\
\midrule
\rowcolor{rowgray} CS2 & SRC     & 412  & 2{,}457  & 5{,}967  & \textbf{95.7}  & 1{,}115  & 10.6 & 12.4 \\
                   CS2 & sCO$_2$ & 340  & 5{,}144  & 15{,}146 & 242.9 & $-$3{,}023 & never & $-$2.3 \\
\rowcolor{rowgray} CS2 & tCO$_2$ & 723  & 4{,}511  & 6{,}238  & 100.0 & \textbf{1{,}710} & 11.4 & 11.6 \\
\midrule
                   CS1 & SRC     & 2639 & 5{,}790  & 2{,}194  & \textbf{35.2}  & \textbf{11{,}415} & 4.2 & 27.8 \\
\rowcolor{rowgray} CS1 & sCO$_2$ & 1988 & 11{,}715 & 5{,}894  & 94.5  & $-$706 & never & 6.2 \\
                   CS1 & tCO$_2$ & 3418 & 11{,}214 & 3{,}281  & 52.6  & 10{,}083 & 7.3 & 17.2 \\
\midrule
\rowcolor{rowgray} CS3 & SRC     & 2273 & 4{,}368  & 1{,}922  & \textbf{30.8}  & \textbf{12{,}406} & 3.2 & 36.2 \\
                   CS3 & sCO$_2$ & 1280 & 7{,}274  & 5{,}684  & 91.1  & 895 & 16.0 & 8.5 \\
\rowcolor{rowgray} CS3 & tCO$_2$ & 2186 & 7{,}733  & 3{,}537  & 56.7  & 7{,}467 & 7.0 & 17.9 \\
\bottomrule
\end{tabularx}
\end{table}

\begin{findingbox}
\textbf{Finding 22 --- The steam cycle gives the lowest cost of electricity at
every site.}
LCOE is \keyresult{AUD\,96, 35 and 31/MWh} for steam at CS2, CS1 and CS3,
against 100, 53 and 57/MWh for tCO$_2$ and 243, 94 and 91/MWh for sCO$_2$.
At the two hotter sites every cycle except sCO$_2$ generates well below the
site electricity price; at the 200\,\textdegree C Pilbara site steam and
tCO$_2$ sit 23--26\,\% below the AUD\,130/MWh price and sCO$_2$ is almost
twice it. The Interim Report's values of AUD\,13--42/MWh for the best cycles
were lower because they omitted heat rejection, mis-sized the heaters and
assumed a 1\,MWe plant regardless of the stream.
\end{findingbox}

\begin{findingbox}
\textbf{Finding 23 --- On NPV the tCO$_2$ cycle is the better investment at
200\,\textdegree C, because it recovers 75\,\% more power from the same
stream; steam is the better investment at 350 and 500\,\textdegree C.}
At CS2 the tCO$_2$ project is worth \keyresult{AUD\,1.71\,M against 1.12\,M}
for steam, despite a 4\,\% higher LCOE, because 723\,kW earns more than
412\,kW. At CS1 steam leads (11.4\,M against 10.1\,M) even though tCO$_2$
recovers 29\,\% more power, because its plant costs 94\,\% more; at CS3
steam leads clearly (12.4\,M against 7.5\,M). The sCO$_2$ project has
negative NPV at CS2 and CS1 and is marginal at CS3.
\end{findingbox}

\begin{findingbox}
\textbf{Finding 24 --- Discounted payback is 3--4 years for steam at the two
hotter sites and 7 years for tCO$_2$; the 200\,\textdegree C site pays back
in 10--11 years for either.}
Steam returns its capital in \keyresult{4.2 and 3.2 years} at CS1 and CS3
with IRRs of 28 and 36\,\%; tCO$_2$ in 7.3 and 7.0 years at 17--18\,\%. At
CS2 both are 10.6--11.4 years at 12\,\% IRR, only five points above the
discount rate. The sCO$_2$ cycle never recovers its capital at CS2 or CS1 and
takes 16 years at CS3. (Simple payback, reported at the interim stage, is
2.8--9.4 years for the cases with positive NPV and understates the time to
recover capital by 1--7 years.)
\end{findingbox}

\section{Efficiency versus Recovery Depth}
\label{sec:rq3_recovery}

\begin{table}[H]
\centering
\caption{Why the economic ranking differs from the efficiency ranking:
recovered power is the product of efficiency and the depth of heat recovery.}
\label{tab:recovery}
\renewcommand{\arraystretch}{1.25}
\footnotesize
\begin{tabularx}{\textwidth}{>{\raggedright\arraybackslash}p{0.9cm}>{\raggedright\arraybackslash}p{1.0cm}
  >{\centering\arraybackslash}X >{\centering\arraybackslash}X >{\centering\arraybackslash}X
  >{\centering\arraybackslash}X >{\centering\arraybackslash}X >{\centering\arraybackslash}X}
\toprule
\rowcolor{tableheader}
\textbf{Site} & \textbf{Cycle} & $\eta_{th}$ rank & Gas out (\textdegree C) & Heat recovered (kW) & Output (kW) & LCOE rank & NPV rank \\
\midrule
\rowcolor{rowgray} CS2 & SRC     & 2 & 113.6 & 2{,}853  & 412  & 1 & 2 \\
                   CS2 & sCO$_2$ & 3 & 110.0 & 2{,}970  & 340  & 3 & 3 \\
\rowcolor{rowgray} CS2 & tCO$_2$ & 1 & 60.2  & 4{,}614  & 723  & 2 & 1 \\
\midrule
                   CS1 & SRC     & 2 & 147.2 & 11{,}153 & 2639 & 1 & 1 \\
\rowcolor{rowgray} CS1 & sCO$_2$ & 3 & 187.5 & 8{,}937  & 1988 & 3 & 3 \\
                   CS1 & tCO$_2$ & 1 & 101.1 & 13{,}689 & 3418 & 2 & 2 \\
\midrule
\rowcolor{rowgray} CS3 & SRC     & 1 & 164.7 & 7{,}376  & 2273 & 1 & 1 \\
                   CS3 & sCO$_2$ & 3 & 289.1 & 4{,}641  & 1280 & 3 & 3 \\
\rowcolor{rowgray} CS3 & tCO$_2$ & 2 & 151.1 & 7{,}675  & 2186 & 2 & 2 \\
\bottomrule
\end{tabularx}
\end{table}

\begin{findingbox}
\textbf{Finding 25 --- In waste heat recovery, depth of recovery matters as
much as efficiency, and the economics follow their product.}
At CS1 the sCO$_2$ cycle is only 1.4 efficiency points below steam
(22.2 against 23.7\,\%) yet produces \keyresult{25\,\% less power}, because
its recuperated heater inlet of 178\,\textdegree C leaves the gas at
188\,\textdegree C while steam cools it to 147\,\textdegree C. At CS2 the
tCO$_2$ cycle is only 1.2 points above steam in efficiency yet produces
75\,\% more power, because it cools the gas to 60\,\textdegree C against
114\,\textdegree C. The useful figure of merit for a WHR cycle is the net
work per kilogram of waste gas (Eq.~\ref{eq:wgas}), and it is on this
quantity, divided by the cost of the plant that delivers it, that the site
rankings turn. A fixed-output comparison hides this entirely: at a common
1\,MWe, LCOE is simply $0.01604 \times$ TEC (AUD/kW) at these economic
parameters, and would reproduce the capital-cost ordering of
Finding~19 at every site.
\end{findingbox}

\section{Monte Carlo Uncertainty and Breakeven Analysis}
\label{sec:rq3_montecarlo}

\begin{table}[H]
\centering
\caption{Monte Carlo results (10,000 samples; capital $\pm 30$\,\%, price
$\pm 40$\,\%, capacity factor 60--95\,\%) and the reduction in heat exchanger
cost (heater plus recuperator) at which a cycle would match the best LCOE at
its site.}
\label{tab:montecarlo}
\renewcommand{\arraystretch}{1.25}
\setlength{\tabcolsep}{3.5pt}
\scriptsize
\begin{tabular}{@{}l l r l l r l@{}}
\toprule
\rowcolor{tableheader}
\textbf{Site} & \textbf{Cycle} & $P(\mathrm{NPV}>0)$ & NPV P10 / P50 / P90 & LCOE P10 / P50 / P90 & Determ.\ NPV & HX cost cut to \\
\rowcolor{tableheader}
 & & (\%) & (kAUD) & (AUD/MWh) & (kAUD) & match best \\
\midrule
\rowcolor{rowgray} CS2 & SRC     & 72.2  & $-$723 / 691 / 2{,}386 & 77 / 104 / 140 & 1{,}115 & --- (best) \\
                   CS2 & sCO$_2$ & 0.5   & $-$5{,}191 / $-$3{,}310 / $-$1{,}423 & 194 / 265 / 355 & $-$3{,}023 & unreachable \\
\rowcolor{rowgray} CS2 & tCO$_2$ & 67.7  & $-$1{,}541 / 966 / 3{,}957 & 80 / 109 / 146 & 1{,}710 & 11\,\% \\
\midrule
                   CS1 & SRC     & 99.8  & 3{,}800 / 9{,}488 / 16{,}555 & 28 / 38 / 51 & 11{,}415 & --- (best) \\
\rowcolor{rowgray} CS1 & sCO$_2$ & 32.9  & $-$7{,}480 / $-$2{,}028 / 3{,}987 & 76 / 103 / 138 & $-$706 & unreachable \\
                   CS1 & tCO$_2$ & 89.7  & $-$81 / 7{,}623 / 17{,}108 & 42 / 57 / 77 & 10{,}083 & 95\,\% \\
\midrule
\rowcolor{rowgray} CS3 & SRC     & 100.0 & 5{,}208 / 10{,}568 / 17{,}247 & 25 / 34 / 45 & 12{,}406 & --- (best) \\
                   CS3 & sCO$_2$ & 48.8  & $-$3{,}745 / $-$87 / 4{,}079 & 73 / 99 / 133 & 895 & unreachable \\
\rowcolor{rowgray} CS3 & tCO$_2$ & 91.3  & 261 / 5{,}699 / 12{,}434 & 45 / 62 / 83 & 7{,}467 & unreachable \\
\bottomrule
\end{tabular}
\end{table}

The Monte Carlo medians lie below the deterministic values because the
capacity-factor distribution is centred on 77.5\,\% rather than the 85\,\%
baseline; the P10--P90 spread is dominated by the $\pm 40$\,\% price range.

\begin{findingbox}
\textbf{Finding 26 --- Steam at 350 and 500\,\textdegree C is robust to every
combination of cost overrun, price fall and reduced running hours; tCO$_2$
at those sites is robust in nine cases out of ten; the 200\,\textdegree C
site is a one-in-three risk for either.}
$P(\mathrm{NPV}>0)$ is \keyresult{99.8 and 100\,\%} for steam at CS1 and CS3
(P10 NPV +AUD\,3.8\,M and +5.2\,M), 89.7 and 91.3\,\% for tCO$_2$, and 72.2
and 67.7\,\% for steam and tCO$_2$ at CS2. The sCO$_2$ cycle is unviable in
practice: 0.5\,\% at CS2, 33\,\% at CS1 and 49\,\% at CS3, with negative
median NPV at all three.
\end{findingbox}

\begin{findingbox}
\textbf{Finding 27 --- Cheaper heat exchangers would make tCO$_2$ competitive
with steam only at the 200\,\textdegree C site, and would not rescue sCO$_2$
anywhere.}
A \keyresult{11\,\% reduction} in the cost of the tCO$_2$ heater and
recuperator would bring its LCOE level with steam at CS2, where its recovery
advantage is largest; at CS1 a 95\,\% reduction would be needed, and at CS3
even free heat exchangers would leave it above steam, because the remaining
turbomachinery and heat-rejection costs are themselves above the whole steam
plant. For sCO$_2$ the target is not reachable at any site. The competitive
window for CO$_2$ cycles in Australian WHR is therefore the low-temperature,
deep-recovery case, and it depends on a cold heat sink (Finding~12).
\end{findingbox}

\section{Summary of RQ3 Results}
\label{sec:rq3_summary}

Sized to what each cycle can actually recover from the site stream and
costed on a consistent basis, the steam Rankine cycle gives the lowest cost
of electricity at all three sites and the highest NPV at the two sites above
300\,\textdegree C, with discounted paybacks of 3--4 years and IRRs of
28--36\,\%. The transcritical CO$_2$ cycle is the better investment at the
200\,\textdegree C Pilbara site, where its deep heat recovery earns it 53\,\%
more NPV than steam, but only if a 14\,\textdegree C heat sink can be
provided at no cost, which Finding~12 shows it cannot. The supercritical
CO$_2$ cycle is not viable at any of the three sites under 2026 conditions.
The central lesson is that the right figure of merit for a WHR cycle is the
work recovered per kilogram of waste gas, and its cost; thermal efficiency
alone predicts neither.
